% Chapter 0: 引言 - 文献概述与背景知识
% do Carmo - 微分几何

\chapter*{引言：文献概述与背景知识}
\addcontentsline{toc}{chapter}{引言：文献概述与背景知识}
\label{chap:Intro}

\section{总述}

本笔记学习的主题是 \textbf{微分几何（differential geometry）}——研究曲线和曲面几何性质的数学分支。这门学科发端于微积分诞生之时，由 Gauss、Riemann 等数学家发展成熟。

本笔记以 Manfredo P. do Carmo 的经典教材《\textbf{Differential Geometry of Curves and Surfaces}》为主线，系统学习曲线与曲面的局部理论。

\section{文献概述}

\subsection*{Differential Geometry of Curves and Surfaces}
\textbf{作者}: Manfredo P. do Carmo \\
\textbf{出版社}: Dover Publications (Revised \& Updated Second Edition, 2016) \\
\textbf{原版}: Prentice-Hall, 1976 \\
\textbf{页数}: 约 500 页

\textbf{内容简介}:
这是微分几何领域最经典的入门教材之一。do Carmo 来自巴西 IMPA（数学与应用数学研究所），他的这本书以清晰、严谨的风格著称。

全书分为五章：
\begin{itemize}
\item \textbf{第1章} 曲线理论 — 局部几何（曲率、挠率、Frenet-Serret 公式）
\item \textbf{第2章} 正则曲面 — 曲面的定义、切平面、第一基本形式
\item \textbf{第3章} Gauss 映射 — 曲面的曲率、第二基本形式
\item \textbf{第4章} 曲面内在几何 — 等距映射、测地线、Gauss-Bonnet 定理
\item \textbf{第5章} 整体微分几何 — 曲面的整体性质
\end{itemize}

\textbf{写作特点}:
\begin{itemize}
\item 每章开头有引言，说明该章内容和学习路线
\item 大量详细的例子和图示
\item 习题丰富，配有提示和答案
\item 强调几何直观与代数方法的结合
\end{itemize}

\textbf{前置要求}:
\begin{itemize}
\item 线性代数（向量、矩阵、行列式）
\item 多变量微积分（偏导数、链式法则、积分）
\item 微分方程（初步了解即可）
\end{itemize}

\section{背景知识安排}

第 \ref{chap:Intro} 章收集学习曲线理论所需的\textbf{纯数学基础}，包括：

\begin{enumerate}
\item \textbf{欧氏空间 $\mathbb{R}^n$} — 向量、范数、内积的基本性质
\item \textbf{向量值函数} — 求导法则、链式法则
\item \textbf{向量积（叉积）} — $\mathbb{R}^3$ 中特有的运算
\item \textbf{行列式与有向体积} — 标量三重积的几何意义
\end{enumerate}

这些内容大部分来自本科线性代数和微积分课程。这里统一整理，便于后续查阅。

\section{学习路径}

根据 do Carmo 在 Preface 中的建议，本笔记采用以下学习顺序：

\begin{enumerate}
\item \textbf{第1步}: 阅读第 \ref{chap:Intro} 章的背景知识
\item \textbf{第2步}: 学习第 \ref{chap:curves} 章（曲线）— 1-2 节到 1-5 节是核心，1-6、1-7 节是拓展
\item \textbf{第3步}: 学习第 \ref{chap:surfaces} 章（曲面）— 2-2、2-3 节是基础，2-4、2-5 节是核心
\item \textbf{第4步}: 深入第 3-4 章（Gauss 映射、内在几何）
\item \textbf{第5步}: 选读第 5 章（整体微分几何）
\end{enumerate}

\section{与其他教材的关系}

微分几何的标准教材还有：
\begin{itemize}
\item do Carmo（本书）— 注重几何直观，例子丰富
\item Klingenberg — 更代数化，强调抽象理论
\item O'Neill — 入门友好，风格接近 do Carmo
\item Morgan — 更加入门，适合本科生
\end{itemize}

do Carmo 的这本书适合作为\textbf{第一本}微分几何教材。

\userannotation{主文献}{do Carmo - Differential Geometry of Curves and Surfaces}
\userannotation{学习顺序}{先曲线，再曲面，最后整体理论}
